\section{Extended torsion and the bordered triangular system}\label{bpI:p1}

This part collects the linear-algebra and first-order geometry of the isoenergetic scheme of HL19
\S6 (pp.~45--51), for the frozen objects $\bar E_H$, $\mathsf T_H$, $\avg{\mathsf T_H}$ and $\mathsf B_h$
of Definitions~\ref{bpI:p0:extended} and~\ref{bpI:p0:bordered}. It proves five things.
\begin{enumerate}[label=(\arabic*)]
\item $\mathsf T_H=\mathsf B_h(I)$ exactly at the trivial torus $\theta\mapsto(\theta,I)$
  (Lemma~\ref{bpI:p1:trivial}); this is why the nondegeneracy hypothesis of I-HEAD is the isoenergetic
  condition of I-G1.
\item The last row of $\mathsf T_H$ is $\omega^\top-(\eta^L)^\top$ exactly, hence $\omega^\top$ on an invariant
  torus (Lemma~\ref{bpI:p1:last-row}); HL19 Rem.~3.7 only states the invariant case, with a sign slip.
\item The bordered triangular system of HL19 Lemma~6.1 with its estimates (Lemma~\ref{bpI:p1:border-solve}),
  and its exact origin as the reduction of the linearized isoenergetic equation in the frame
  (Lemma~\ref{bpI:p1:reduced}). \emph{This is the only genuinely new analytic lemma}; HL19 print only
  ``analogous to Lemma~5.1'' for its proof, so the full proof is supplied.
\item The approximate energy identities (Lemma~\ref{bpI:p1:energy-approx}): $D(H\circ K)=DH(K)DK$,
  $\partial_\omega(H\circ K)=-DH(K)E$, the compatibility $\avg{\eta^N}=0$, and the size of the energy defect.
\item Energy conservation of an invariant torus (Lemma~\ref{bpI:p1:energy-conservation}): $E=0$ and
  $\avg{H\circ K}=E_0$ give $H\circ K=E_0$ pointwise, the pointwise energy conclusion (O3) of I-G1.
\end{enumerate}
Everything about the frame $L=DK$, $N$, $P$, $Q$, the symplectic defects and the cohomological equation is
\emph{cited} from track~K's Part~2 (\Kref{bp:p2:frame-identities}, \Kref{bp:p2:exactness},
\Kref{bp:p2:geometry}, \Kref{bp:p2:cohomological}, \Kref{bp:p2:triangular}, \dots) and Part~1, not
re-proved: the canonical Case~III objects of the frozen definitions are those of
\Kref{bp:p2:objects} with $h$ replaced by $H$.

%----------------------------------------------------------------------------------------------------
\subsection{Setting, sign dictionary, external facts}\label{bpI:p1:sec-setting}

\begin{assumption}[Standing setting of Part~\ref{bpI:p1}]\label{bpI:p1:setting}
Let $n\ge2$ and $\tau\ge0$ (I-G1 has $\tau\ge n-1$; only $\tau\ge0$ is used here), $M>0$. Let
$D\subset\C^n\times\C^n$ be open with $(\theta+k,I)\in D$ for $(\theta,I)\in D$, $k\in\Z^n$ (I1); let
$H:\C^n\times\C^n\to\C$ be holomorphic on $D$, $1$-periodic in $\theta$, real at real phase points of $D$
(I2, I3), with $\norm{DH},\norm{X_H},\norm{DX_H},\norm{D^2X_H}\le M$ on $D$ (I4). Let $\gamma>0$ and
$\omega\in\R^n\cap\Dio_{\gamma,\tau}$. Let $r>0$ and let $K:\C^n\to\C^n\times\C^n$ be holomorphic, primary and real
on reals on $T_r$, with $K(T_r)\subset D$ (I5, I6). Where the frame is used, $\det G_K\ne0$ on $T_r$.
\end{assumption}

This is exactly Assumption~\Kref{bp:p2:standing} with $h:=H$, $\mathfrak B:=D$, together with the bound
$\norm{DH}\le M$, which is used only in Lemma~\ref{bpI:p1:energy-approx}(e) and
Lemma~\ref{bpI:p1:extended-basic}(b). All objects $L,G,B,N,P,X_N,T,E,A,\Omega_K,E_{\mathrm{sym}},Q,\eta$ and the
averages are those of \Kref{bp:p2:objects}, with $h:=H$; \Kref{bp:p2:frozen-agreement} identifies them with the
frozen \Lean{gramMatrix}, \Lean{normalFrame}, \Lean{torsion}, \Lean{invarianceError}. The norms are those of
Part~\ref{bpI:p0}: max norms, row-sum operator norms $\abs M=\max_i\sum_j\abs{M_{ij}}$ (so a $1\times n$ matrix has
the norm $\sum_j\abs{M_{1j}}$, which is the norm of $DH(z)$ in (I4)). Write $\Gamma_\delta:=c_R\gamma^{-1}\delta^{-\varkappa}$,
$\varkappa=\tau+n$, as in \Kref{bp:p2:constants}.

\begin{remark}[Rays: $s\omega^*$ is Diophantine with the constant of $\omega^*$]\label{bpI:p1:ray}
In I-G1, $\omega=s\omega^*$ with $s>1$ and $\omega^*\in\Dio_{\gamma,\tau}$. For $k\ne0$,
$\abs{k\cdot s\omega^*}=s\abs{k\cdot\omega^*}\ge\gamma/\bigl(\sum_i\abs{k_i}\bigr)^\tau$, so
$\omega\in\Dio_{\gamma,\tau}$ with the \emph{same} $\gamma$. Hence every result of K's Part~1 and Part~2 that
requires $\omega\in\Dio_{\gamma,\tau}$ applies to $\omega=s\omega^*$ with $\Gamma_\delta$ as above.
\end{remark}

\begin{remark}[Lean totality for this part]\label{bpI:p1:totality}
Following Remark~\ref{bpI:p0:totality}: every statement below that uses $N$, $T$, $\mathsf T_H$ assumes
$\det G_K\ne0$ on $T_r$ first, so that $B=G^{-1}$ is the true inverse, holomorphic
(\Kref{bp:p1:holomorphic-inverse}); every use of $\avg{\mathsf T_H}$ is preceded by continuity of
$\mathsf T_H$ on $\R^n$ (Lemma~\ref{bpI:p1:extended-basic}(a)), hence integrability
(\Kref{bp:p1:ext-integral}); every derivative of $H$, $X_H$ is taken at a point of the open set $D$ where $H$ is
holomorphic, so \Lean{fderiv} is the genuine derivative; the frozen $\mathsf T_H$ entry
$\Lean{fderiv}\ \C\ H\ (K\theta)\ (N_{\cdot b})$ is then the matrix product $DH(K(\theta))N_{\cdot b}(\theta)$.
\end{remark}

\begin{remark}[Sign dictionary HL19 $\leftrightarrow$ this blueprint]\label{bpI:p1:signs}
HL19 use $\mathfrak L_\omega=-\sum_i\omega_i\partial_{\theta_i}=-\partial_\omega$ (4.1), p.~20, and the
right operator $\mathfrak R_\omega$ of (4.2), $\widehat{(\mathfrak R_\omega v)}_k=-\hat v_k/(2\pi ik\cdot\omega)$. Track~K's operator
$\mathcal R$ (\Kref{bp:p2:cohomological}) solves $\partial_\omega u=g$, $\avg u=0$, i.e.
$\widehat{(\mathcal Rg)}_k=\hat g_k/(2\pi ik\cdot\omega)$. Thus
\[
\mathfrak L_\omega=-\partial_\omega,\qquad\mathfrak R_\omega=-\mathcal R\quad\text{(on the data of (4.1))}.
\]
With the canonical $\Omega=\Omega_0$ and $\Omega X_h=(Dh)^\top$ (Rem.~2.1, p.~6; our
\Kref{bp:p2:hamiltonian}(d)), HL19's invariance error is $E=X_h(K)-DK\,\omega$ (2.2), p.~7, which p.~45 writes as
$X_h(K)+\mathfrak L_\omega K$. The linearized system (6.1), p.~45, with $\bar\omega=\omega+\Delta\omega$,
$\Delta\omega=-\omega\xi^\omega$ (6.3), is therefore
\[
A\,\Delta K-\partial_\omega\Delta K-DK\,\Delta\omega=A\,\Delta K-\partial_\omega\Delta K+L\omega\,\xi^\omega=-E,
\qquad\avg{DH(K)\Delta K}=-E^\omega,
\]
$A=DX_H\circ K$, $L=DK$, $E^\omega=\avg{H\circ K}-E_0$, and $\bar\omega=(1-\xi^\omega)\omega$ as on p.~50. The system (6.7) is
\[
T\xi^N-\partial_\omega\xi^L+\omega\,\xi^\omega=\eta^L,\qquad-\partial_\omega\xi^N=\eta^N,\qquad\avg{Y\xi^N}=\eta^\omega,
\]
and the solution (6.8)--(6.10) reads $\xi^N=\xi^N_0-\mathcal R\eta^N$, $\xi^L=\xi^L_0+\mathcal R\bigl(T\xi^N+\omega\xi^\omega-\eta^L\bigr)$,
with (6.10) unchanged because $\mathfrak R_\omega=-\mathcal R$. Lemma~\ref{bpI:p1:reduced} verifies that this is the
equation obtained from (6.1) in the frame, signs included. Here and below $Y(\theta):=DH(K(\theta))N(\theta)\in\C^{1\times n}$ is HL19's
$H$ of (6.14), renamed to avoid a clash with the Hamiltonian.
\end{remark}

\paragraph{External facts used in this part.} Besides \Kref{bp:p1:ext-integral} (continuous functions on a compact
box are integrable, integrals are linear and bounded by the supremum), the proofs use only:
\begin{enumerate}[label=(X\arabic*),ref=(X\arabic*)]
\item\label{bpI:p1:ext-chain} \emph{Chain and product rules} for Fr\'echet derivatives of holomorphic maps, and
  linearity of the derivative (Mathlib \nolinkurl{HasFDerivAt.comp}, \nolinkurl{fderiv_add}, \nolinkurl{fderiv_const}).
\item\label{bpI:p1:ext-real-direction} \emph{Complex derivative in a real direction.} If $\varphi$ is
  holomorphic on an open $U\subset\C^n$, $x\in U\cap\R^n$ and $v\in\R^n$, then $\varphi|_{U\cap\R^n}$ has the
  real directional derivative $D\varphi(x)v$ at $x$; if moreover $\varphi$ is real on $U\cap\R^n$, then
  $D\varphi(x)v\in\R$ (it is a limit of real difference quotients). See \Kref{bp:p1:hamiltonian-real},
  step \Kref{step:p1:hr-derivative}.
\end{enumerate}

%----------------------------------------------------------------------------------------------------
\subsection{The extended torsion}\label{bpI:p1:sec-extended}

\begin{lemma}[Regularity, periodicity, realness and size of the extended torsion]\label{bpI:p1:extended-basic}
Assume~\ref{bpI:p1:setting} with $\det G_K\ne0$ on $T_r$, and put $Y:=DH(K)N$ on $T_r$. Then:
\begin{enumerate}[label=(\alph*)]
\item $Y$ and the entries of $\mathsf T_H$ (built from $(H,K,\omega)$) are holomorphic and $1$-periodic on $T_r$ and
  real at real $\theta$. In particular $\mathsf T_H$ is continuous on $\R^n$ and $\avg{\mathsf T_H}$ is defined, real, and equals
  the block matrix
  \[
  \avg{\mathsf T_H}=\begin{pmatrix}\avg T&\omega\\ \avg Y&0\end{pmatrix}
  \]
  with $\avg T=\avg{T_{H,K,\omega}}$ the torsion average of \Kref{bp:p2:objects}.
\item If $\abs{DK}\le\sigma_K'$ and $\abs B\le\sigma_B'$ on $T_r$, then $\abs Y\le M\sigma_K'\sigma_B'$ on $T_r$, and
  on any set where $\abs T\le C_T$,
  $\abs{\mathsf T_H}\le\max\{C_T+\abs\omega,\ M\sigma_K'\sigma_B'\}$ in the row-sum norm of $\C^{(n+1)\times(n+1)}$.
\end{enumerate}
\frozen{KAM.Iso.extendedTorsion, KAM.Iso.extendedTorsionAverage}
\source{HL19 (3.13), p.~19 ($c=H$, $d=n$); HL19 (6.14), p.~47 for $Y$. Estimates: ours.}
\end{lemma}

\begin{proof}
\proofstep{Holomorphy}\label{stepI:p1:eb-holo}
\sloppy $N$ is holomorphic on $T_r$ (\Kref{step:p2:fb-holo}, which uses only $K(T_r)\subset D$, $\det G_K\ne0$ and the holomorphic inverse \Kref{bp:p1:holomorphic-inverse}; no bound and no torsion condition enters), and $DH$ is holomorphic on $D$
(\Kref{bp:p1:holomorphic-regularity}), so $Y=DH\circ K\cdot N$ is a product of holomorphic functions
(\ref{bpI:p1:ext-chain}). The entries of $T$ and the constant block $\omega$ are holomorphic by the same step. Hence every entry of $\mathsf T_H$ is holomorphic, in particular continuous on $\R^n\subset T_r$.

\proofstep{Periodicity and realness}\label{stepI:p1:eb-period}
Differentiating $H(\theta+k,I)=H(\theta,I)$ gives $DH(z+(k,0))=DH(z)$ on $D$, and $K(\theta+k)=K(\theta)+(k,0)$;
so $DH(K(\theta+k))=DH(K(\theta))$. $N$ and $T$ are $1$-periodic by \Kref{bp:p2:periodicity}. At real $\theta$,
$K(\theta)$ is a real phase point of $D$, where $DH$ is real (the argument of \Kref{bp:p2:hamiltonian}(c),
step \Kref{step:p2:ham-real}, applies to $H$, real at real points), and $N$, $T$ are real
(\Kref{bp:p2:periodicity}). Hence $Y$ is real at real points.

\proofstep{The average}\label{stepI:p1:eb-average}
\sloppy By Step~\ref{stepI:p1:eb-holo}, each entry is integrable on $[0,1]^n$
(\Kref{bp:p1:ext-integral}), so the Bochner
integral of the frozen $\avg{\mathsf T_H}$ is the entrywise integral, real because the integrand is
real at real $\theta$; the entries of the border are $\avg\omega=\omega$ (the cube has volume one) and $\avg 0=0$.
(Lean totality, Remark~\ref{bpI:p1:totality}.)

\proofstep{Bounds}\label{stepI:p1:eb-bounds}
$\abs Y\le\abs{DH(K)}\abs N\le M\,\abs{\Omega_0}\abs L\abs B\le M\sigma_K'\sigma_B'$, by \Kref{bp:p2:matrix}(b), (d).
The last column of $\mathsf T_H$ is the $n\times1$ matrix $\omega$, of norm $\abs\omega$; by
\Kref{bp:p2:matrix}(c) the $(n+1)\times(n+1)$ norm is at most the maximum of the two block-row bounds
$C_T+\abs\omega$ and $\abs Y$.
\end{proof}

\begin{lemma}[$\mathsf T_H=\mathsf B_h$ at the trivial torus]\label{bpI:p1:trivial}
Let $U\subset\C^n$ be open and $h:U\to\C$ holomorphic with $h(J)\in\R$ for every real $J\in U$. Let
$D=\C^n\times U$, $H(\theta,J):=h(J)$ on $D$, $I\in U\cap\R^n$ and $K(\theta)=(\theta,I)$ ($\theta\in\C^n$). Put
$\omega:=\nabla h(I)\in\R^n$ (\Lean{frequencyMap}). Then $K$ is holomorphic, primary and real on reals, $K(\C^n)\subset D$,
$G_K=I_n$, $E_{H,K,\omega}\equiv0$, and for every $\theta\in\C^n$
\[
\mathsf T_H(\theta)=\mathsf B_h(I)\quad\text{(the real matrix, viewed in }\C^{(n+1)\times(n+1)}),\qquad
\avg{\mathsf T_H}=\mathsf B_h(I),
\]
the frozen \Lean{borderedHessian}. In particular $\det\avg{\mathsf T_H}\ne0$ iff $\det\mathsf B_h(I)\ne0$, and the (total, Lean) inverses of the two matrices correspond under the cast, so
$\abs{\avg{\mathsf T_H}^{-1}}=\abs{\mathsf B_h(I)^{-1}}$ in both cases.
\frozen{KAM.Iso.extendedTorsion, KAM.Iso.extendedTorsionAverage, KAM.Iso.borderedHessian}
\source{ours; HL19 Rem.~3.7, p.~19, and Chi05 p.~17 (determinant $3/L^4$ in the Kepler case). Reading remark:
Remark~\ref{bpI:p0:head-remarks}(b).}
\end{lemma}

\begin{proof}
\proofstep{The torus}\label{stepI:p1:tr-torus}
$K(\theta+k)=(\theta+k,I)=K(\theta)+(k,0)$, $K$ is affine, hence holomorphic, and real on real $\theta$ because $I$
is real. $DK(\theta)e_i=(e_i,0)$, so $G_K=(e_i\cdot e_j)=I_n$ (using the bilinear dot of
Definition~\ref{bpI:p0:stripdata}), $\det G_K=1$ and $B=I_n$.

\proofstep{The normal frame}\label{stepI:p1:tr-frame}
$N_{\cdot j}=\sum_k\delta_{kj}\Omega_0(e_k,0)=\Omega_0(e_j,0)=(0,e_j)$, a constant function of $\theta$, so
$D_\theta N_{\cdot j}=0$.

\proofstep{The vector field}\label{stepI:p1:tr-field}
Since $H(\theta,J)=h(J)$, $X_H(\theta,J)=\bigl((\partial_ah(J))_a,\,0\bigr)$, so $E_{H,K,\omega}(\theta)=X_H(K\theta)-DK\,\omega=
(\nabla^{\C}h(I)-\omega,0)$, where $\nabla^\C h(I)=(\partial_ah(I))_a$. By~\ref{bpI:p1:ext-real-direction} applied to
$\varphi=h$ (real on $U\cap\R^n$), $\partial_ah(J)\in\R$ for real $J$, and the real function
$J\mapsto\operatorname{Re}h(J)=h(J)$ on $U\cap\R^n$ has partial derivative $\partial_ah(J)$ in the direction $e_a$; by the
definition of \Lean{frequencyMap}, $\omega_a=\partial_ah(I)$. Hence $E\equiv0$, and more generally
$\nabla h(J)_a=\partial_ah(J)$ for real $J\in U$.

\proofstep{The torsion}\label{stepI:p1:tr-torsion}
$DX_H(\theta,I)(0,e_j)=\bigl((\partial_j\partial_ah(I))_a,0\bigr)$, and $\Omega_0(w,0)=(0,w)$. Hence
$T(\theta)_{ij}=N_{\cdot i}\cdot\Omega_0\bigl(DX_H N_{\cdot j}-0\bigr)=(0,e_i)\cdot(0,(\partial_j\partial_ah(I))_a)=\partial_j\partial_ih(I)$.

\proofstep{The border}\label{stepI:p1:tr-border}
The last column is $\omega=\nabla h(I)$, entries $\omega_a=\partial_ah(I)$. The last row is
$DH(K\theta)N_{\cdot b}=DH(\theta,I)(0,e_b)=\partial_bh(I)=\omega_b$. The corner is $0$.

\proofstep{Comparison with the bordered Hessian}\label{stepI:p1:tr-compare}
The $(a,b)$ entry of $\mathsf B_h(I)$ is $D(\nabla h)(I)e_b$ read at $a$, the real derivative in direction $e_b$ of
$J\mapsto\nabla h(J)_a=\partial_ah(J)$ on $U\cap\R^n$ (Step~\ref{stepI:p1:tr-field}); since $\partial_ah$ is holomorphic,
this is $\partial_b\partial_ah(I)$ by~\ref{bpI:p1:ext-real-direction}, a real number. So the upper left blocks agree
(and are symmetric by \Kref{bp:p1:holomorphic-regularity}); the border of $\mathsf B_h(I)$ is $\nabla h(I)$ in the
last column and row, as found above. Thus $\mathsf T_H(\theta)=\mathsf B_h(I)$ for all $\theta$; the constant has
average equal to itself since $[0,1]^n$ has volume one. The statements on the determinant and the inverse follow by cases. If $\det\mathsf B_h(I)\ne0$, a real matrix and its image in $\C^{(n+1)\times(n+1)}$ have the same determinant, and inverses correspond
(\Kref{bp:p1:real-inverse}); the row-sum norm is unchanged. If $\det\mathsf B_h(I)=0$, the cast has determinant $0$ and both total inverses are the zero matrix, so the inverse norms are both $0$ (the statement is then vacuous for (I8), which requires $\det\ne0$).
\end{proof}

\begin{remark}[Use]\label{bpI:p1:trivial-use}
For $H=h+f$ with $f$ small, the torus $K(\theta)=(\theta,I)$ with $\omega=\nabla h(I)$ (or $s\omega^*$ with a
nearby $I$, see Part~3) has $\avg{\mathsf T_H}=\mathsf B_h(I)+O(\norm f)$ by Lemma~\ref{bpI:p1:last-row}
below and the Cauchy estimates; this is how (I8) of I-G1 is verified from $\det\mathsf B_h\ne0$, with
$\sigma_T$ from \Kref{bp:p1:neumann}.
\end{remark}

\begin{lemma}[The energy row of the extended torsion]\label{bpI:p1:last-row}
Assume~\ref{bpI:p1:setting} with $\det G_K\ne0$ on $T_r$ and let $E=E_{H,K,\omega}$, $\eta^L=-N^\top\Omega_0E$ as in
\Kref{bp:p2:objects}. Then on $T_r$:
\begin{enumerate}[label=(\alph*)]
\item $Y=DH(K)N=\omega^\top-(\eta^L)^\top$. In particular $\abs{Y-\omega^\top}=\sum_b\abs{\eta^L_b}\le n\abs{\eta^L}$ (row-sum norm of a $1\times n$ matrix) and, entrywise,
  $Y_b=\omega_b+N_{\cdot b}^\top\Omega_0E$.
\item If $E=0$ on $T_r$ (invariant torus), then $Y=\omega^\top$ on $T_r$, so
  \[
  \mathsf T_H(\theta)=\begin{pmatrix}T(\theta)&\omega\\ \omega^\top&0\end{pmatrix}\quad(\theta\in T_r),\qquad
  \avg{\mathsf T_H}=\begin{pmatrix}\avg T&\omega\\ \omega^\top&0\end{pmatrix}.
  \]
\item In general $\avg{\mathsf T_H}=\mathsf T^{\mathrm{inv}}-\begin{pmatrix}0&0\\ \avg{\eta^L}^\top&0\end{pmatrix}$, where
  $\mathsf T^{\mathrm{inv}}:=\bigl(\begin{smallmatrix}\avg T&\omega\\ \omega^\top&0\end{smallmatrix}\bigr)$ is built from the same $T$; and if
  $\abs E\le e$ on $T_r$, $\abs{L}\le\sigma_K'$, $\abs B\le\sigma_B'$, then
  $\abs{\avg{\mathsf T_H}-\mathsf T^{\mathrm{inv}}}\le2n^2\sigma_K'\sigma_B'\,e$.
\end{enumerate}
\frozen{KAM.Iso.extendedTorsion}
\source{HL19 Rem.~3.7, p.~19 (part (b), (3.17)), proved with the sign corrected (Remark~\ref{bpI:p1:errata}, E3);
(a), (c) are ours: HL19 (3.13) defines the last row as $Dc(K)N$ without simplifying it off invariance.}
\end{lemma}

\begin{proof}
\proofstep{The gradient identity}\label{stepI:p1:lr-gradient}
By \Kref{bp:p2:hamiltonian}(d) (equivalently \Kref{bp:p1:hamiltonian-real}), $\Omega_0X_H=(DH)^\top$ on $D$, so
$DH(z)=(\Omega_0X_H(z))^\top=X_H(z)^\top\Omega_0^\top=-X_H(z)^\top\Omega_0$.

\proofstep{Insert the error}\label{stepI:p1:lr-insert}
At $z=K(\theta)$, $X_H(K)=L\omega+E$ by the definition of $E$. Hence
\[
Y=DH(K)N=-(L\omega+E)^\top\Omega_0N=-\omega^\top L^\top\Omega_0N-E^\top\Omega_0N=\omega^\top-E^\top\Omega_0N,
\]
using $L^\top\Omega_0N=-I_n$ (\Kref{bp:p2:frame-identities}(a)). The $b$-th entry of the row $E^\top\Omega_0N$ is the scalar
$E^\top\Omega_0N_{\cdot b}=N_{\cdot b}^\top\Omega_0^\top E=-N_{\cdot b}^\top\Omega_0E=(\eta^L)_b$. This is (a) and the entrywise form. The row-sum norm of
the $1\times n$ matrix $(\eta^L)^\top$ is $\sum_b\abs{\eta^L_b}\le n\abs{\eta^L}$, so $\abs{Y-\omega^\top}=\sum_b\abs{\eta^L_b}\le n\abs{\eta^L}$ (the factor $n$ cannot be dropped: for $H=(\omega-(1,\dots,1))\cdot I$, $K=(\theta,0)$ one has $\eta^L=(1,\dots,1)$ and $\abs{Y-\omega^\top}=n$).

\proofstep{Invariant torus and average}\label{stepI:p1:lr-average}
(b) is (a) with $\eta^L=0$ (since $E=0$), together with the definition of $\mathsf T_H$
(Definition~\ref{bpI:p0:extended}) and Lemma~\ref{bpI:p1:extended-basic}(a). For (c): averaging (a) and using linearity,
$\avg Y=\omega^\top-\avg{\eta^L}^\top$, so the two matrices differ only in the last row, which has norm
$\le n\abs{\avg{\eta^L}}\le n\sup_{\R^n}\abs{\eta^L}$ (\Kref{bp:p2:averages}(a)), and
$\abs{\eta^L}\le\abs{N^\top}\abs{\Omega_0}\abs E\le2n\,\abs N\abs E\le2n\,\sigma_K'\sigma_B'\,e$ ($\abs{N^\top}\le2n\abs N$)
(\Kref{bp:p2:matrix}(b), (d), with $\abs N\le\abs L\abs B$). Hence the bound $2n^2\sigma_K'\sigma_B'e$.
\end{proof}

%----------------------------------------------------------------------------------------------------
\subsection{Energy identities}\label{bpI:p1:sec-energy}

\begin{lemma}[Approximate energy identities]\label{bpI:p1:energy-approx}
Assume~\ref{bpI:p1:setting} and let $E=E_{H,K,\omega}$, $E^{\mathrm{en}}:=H\circ K-\avg{H\circ K}$ on $T_r$. Then:
\begin{enumerate}[label=(\alph*)]
\item $H\circ K$ is holomorphic, $1$-periodic on $T_r$, real on reals, and $D(H\circ K)=DH(K)\,DK$; the average
  $\avg{H\circ K}=\bar E_H(K)$ is the frozen \Lean{energyAverage}, real.
\item $\partial_\omega(H\circ K)=D(H\circ K)\omega=-DH(K)\,E$ on $T_r$.
\item $L^\top\Omega_0E=(D(H\circ K))^\top-\Omega_K\omega$ on $T_r$. If $\det G_K\ne0$, then $\eta^N=L^\top\Omega_0E$.
\item $\avg{D(H\circ K)}=0$, $\avg{\Omega_K}=0$ and $\avg{\eta^N}=0$ (the compatibility condition of Lemma~\ref{bpI:p1:border-solve}).
\item If $\abs E\le\varepsilon$ on $T_r$ and $0<\delta<r$, then $\abs{E^{\mathrm{en}}}\le M\Gamma_\delta\varepsilon$ on $T_{r-\delta}$, and for
  $0<\delta'<r-\delta$, $\abs{D(H\circ K)}\le nM\Gamma_\delta\varepsilon/\delta'$ on $T_{r-\delta-\delta'}$.
\item (Energy defect.) Let $\det G_K\ne0$ on $T_r$, let $\xi=(\xi^L,\xi^N)$ be holomorphic on a strip $T_{r'}$ with $0<r'\le r$ (so $\xi$ is continuous on the real cube), and
  $\Delta K=P\xi$. Then on $T_{r'}$
  \[
  DH(K)\,\Delta K=Y\xi^N+D(H\circ K)\,\xi^L=Y\xi^N+(\eta^N)^\top\xi^L-\omega^\top\Omega_K\,\xi^L .
  \]
  Consequently, if $\avg{Y\xi^N}=\eta^\omega$, then $\avg{DH(K)\Delta K}=\eta^\omega+E^\omega_{\mathrm{lin}}$ with
  $E^\omega_{\mathrm{lin}}:=\avg{D(H\circ K)\,\xi^L}$, and
  $\abs{E^\omega_{\mathrm{lin}}}\le\sup_{[0,1]^n}\abs{D(H\circ K)}\,\sup_{[0,1]^n}\abs{\xi^L}\le n\bigl(\sup_{[0,1]^n}\abs{\eta^N}+\abs\omega\sup_{[0,1]^n}\abs{\Omega_K}\bigr)\sup_{[0,1]^n}\abs{\xi^L}$, all suprema finite by continuity on the compact cube.
\end{enumerate}
\frozen{KAM.Iso.energyAverage}
\source{HL19 Lemma~4.2, (4.3)--(4.5), p.~21, and (4.9) ($\mathfrak L_\omega(c\circ K)=Dc(K)E$, equivalent to (b) with
$\mathfrak L_\omega=-\partial_\omega$); p.~46 ($Dc(K)L=D(c\circ K)$, $E^\omega_{\mathrm{lin}}$ of p.~49); (6.11)--(6.14), p.~47;
compatibility $\avg{\eta^N}=0$ from Lemma~4.6 as recalled on p.~47.}
\end{lemma}

\begin{proof}
\proofstep{Regularity and chain rule}\label{stepI:p1:ea-chain}
$K$ is holomorphic on $T_r$ with values in $D$, and $H$ is holomorphic on $D$, so $H\circ K$ is holomorphic with
$D(H\circ K)=DH(K)DK$ (\ref{bpI:p1:ext-chain}). It is $1$-periodic since $H(K(\theta+k))=H(K(\theta)+(k,0))=H(K(\theta))$,
and real at real $\theta$ because $K(\theta)$ is a real phase point of $D$ (I3). The average is
$\int_{[0,1]^n}H(K(\theta))d\theta=\bar E_H(K)$ by definition, integrable by \Kref{bp:p2:averages}(a). This is (a).

\proofstep{The derivative along $\omega$}\label{stepI:p1:ea-omega}
$D(H\circ K)\omega=DH(K)L\omega=DH(K)\bigl(X_H(K)-E\bigr)$. By \Kref{bp:p2:hamiltonian}(d), $DH=(\Omega_0X_H)^\top$,
so $DH(z)X_H(z)=X_H^\top\Omega_0^\top X_H=-X_H^\top\Omega_0X_H=0$ ($\Omega_0$ is antisymmetric). Hence $D(H\circ K)\omega=-DH(K)E$, which is (b).

\proofstep{The gradient identity}\label{stepI:p1:ea-gradient}
This repeats the gradient step of \Kref{bp:p2:exactness}: $L^\top\Omega_0X_H(K)=(DH(K)L)^\top=(D(H\circ K))^\top$,
hence $L^\top\Omega_0E=(D(H\circ K))^\top-L^\top\Omega_0L\omega=(D(H\circ K))^\top-\Omega_K\omega$. The identification $\eta^N=L^\top\Omega_0E$ is
the definition of $\eta$ in \Kref{bp:p2:objects}. This is (c).

\proofstep{Averages}\label{stepI:p1:ea-avg}
$D(H\circ K)$ has the entries $\partial_{\theta_\ell}(H\circ K)$ of a holomorphic periodic function, so the average vanishes
(\Kref{bp:p2:averages}(b)); $\avg{\Omega_K}=0$ and $\avg{\eta^N}=0$ are \Kref{bp:p2:exactness}. (Also directly:
$\avg{\eta^N}=\avg{D(H\circ K)}^\top-\avg{\Omega_K}\omega=0$ by (c).)

\proofstep{The size of the energy fluctuation}\label{stepI:p1:ea-size}
$E^{\mathrm{en}}$ is holomorphic and $1$-periodic on $T_r$, has zero average, and by (b)
$\abs{\partial_\omega E^{\mathrm{en}}}=\abs{DH(K)E}\le\abs{DH(K)}\abs E\le M\varepsilon$ on $T_r$ (here (I4), $K(T_r)\subset D$, is used).
\Kref{bp:p2:cohomological}(b), on $T_r$ with loss $\delta$, gives $\abs{E^{\mathrm{en}}}\le\Gamma_\delta M\varepsilon$ on $T_{r-\delta}$. The
Cauchy estimate \Kref{bp:p1:cauchy-strip} for the bounded holomorphic function $E^{\mathrm{en}}$ on $T_{r-\delta}$ with loss $\delta'$ gives
$\abs{DE^{\mathrm{en}}}=\abs{D(H\circ K)}\le n\Gamma_\delta M\varepsilon/\delta'$ on $T_{r-\delta-\delta'}$. This is (e).

\proofstep{The energy defect}\label{stepI:p1:ea-defect}
$P=(L\ N)$, so $\Delta K=L\xi^L+N\xi^N$ and $DH(K)\Delta K=DH(K)L\xi^L+DH(K)N\xi^N=D(H\circ K)\xi^L+Y\xi^N$ by (a).
By (c), $D(H\circ K)=(\eta^N+\Omega_K\omega)^\top=(\eta^N)^\top+\omega^\top\Omega_K^\top=(\eta^N)^\top-\omega^\top\Omega_K$ since
$\Omega_K^\top=-\Omega_K$ ($\Omega_K=L^\top\Omega_0L$, $\Omega_0^\top=-\Omega_0$). Since $r'>0$, every term is continuous on $[0,1]^n\subset T_{r'}$, hence integrable (\Kref{bp:p1:ext-integral}), and the average of the sum splits linearly.
For the bound: $\abs{\avg{D(H\circ K)\xi^L}}\le\sup_{[0,1]^n}\abs{D(H\circ K)\xi^L}$ (\Kref{bp:p1:ext-integral}, volume one), and
$\abs{D(H\circ K)}\le\abs{(\eta^N)^\top}+\abs{\omega^\top\Omega_K}\le n\abs{\eta^N}+n\abs\omega\abs{\Omega_K}$, with
\Kref{bp:p2:matrix}(a) for transposes (a $1\times n$ row $v^\top$ has $\abs{v^\top}=\sum\abs{v_i}\le n\abs v$).
\end{proof}

\begin{lemma}[Energy conservation of an invariant torus]\label{bpI:p1:energy-conservation}
Assume~\ref{bpI:p1:setting} on a strip $T_r$ (with $\omega\in\Dio_{\gamma,\tau}$) and $E_{H,K,\omega}(\theta)=0$ for all $\theta\in T_r$.
Then:
\begin{enumerate}[label=(\alph*)]
\item $H\circ K\equiv\bar E_H(K)$ is constant on $T_r$.
\item If $E_0\in\R$ (or $\C$) and $\bar E_H(K)=E_0$, then $H(K(\theta))=E_0$ for every $\theta\in T_r$.
\item Quantitatively, if $E_{H,K,\omega}(\theta)=0$ on $T_r$ then $\sup_{T_r}\abs{H\circ K-E_0}=\abs{\bar E_H(K)-E_0}$.
\end{enumerate}
\frozen{KAM.Iso.energyAverage}
\source{HL19 Rem.~2.2, p.~7 (``quasi-periodicity implies $h(K(\theta))=\avg{h\circ K}$''); here proved. This is the pointwise
energy conclusion (O3) of I-G1 (decision D-I-ENERGY).}
\end{lemma}

\begin{proof}
\proofstep{The derivative vanishes along $\omega$}\label{stepI:p1:ec-derivative}
Let $u:=H\circ K-\avg{H\circ K}$, holomorphic and $1$-periodic on $T_r$ with $\avg u=0$ (Lemma~\ref{bpI:p1:energy-approx}(a)).
By Lemma~\ref{bpI:p1:energy-approx}(b), $Du\,\omega=-DH(K)E=0$ on $T_r$.

\proofstep{Uniqueness for the cohomological equation}\label{stepI:p1:ec-unique}
For every $0<\delta<r$, $u|_{T_{r-\delta}}$ is holomorphic, $1$-periodic, of zero average, and solves $Du\,\omega=g$ with
$g=0$ (holomorphic, bounded, periodic, $\avg g=0$ on $T_r$). By the uniqueness clause of \Kref{bp:p1:cohomological}
(no boundedness of the solution assumed), $u=0$ on $T_{r-\delta}$. As $\delta\in(0,r)$ is arbitrary, $u=0$ on $T_r=\bigcup_\delta T_{r-\delta}$,
which is (a).

\proofstep{Pointwise energy}\label{stepI:p1:ec-pointwise}
(b) and (c) are (a): $H\circ K\equiv\bar E_H(K)$, so $\abs{H(K(\theta))-E_0}=\abs{\bar E_H(K)-E_0}$ for every $\theta$.
\end{proof}

\begin{remark}[Where (O3) of I-G1 comes from]\label{bpI:p1:o3}
In the proof of I-G1 (Part~2), $E_{H,K_\infty,s_\infty\omega^*}=0$ on $T_{\rho_\infty}$ and $\bar E_H(K_\infty)=E_0$: the second follows
from $\abs{\bar E_{H}(K_j)-E_0}\to0$ along the iteration and the uniform convergence $K_j\to K_\infty$ (continuity of $H$ and of the
integral on the compact cube), which is part of Part~2; then Lemma~\ref{bpI:p1:energy-conservation}(b), with $\omega=s_\infty\omega^*\in\Dio_{\gamma,\tau}$
(Remark~\ref{bpI:p1:ray}), gives $H(K_\infty)=E_0$ on $T_{\rho_\infty}$. Neither the Gram determinant of $K_\infty$ nor the frame
is used (Remark~\ref{bpI:p0:totality}(T2)).
\end{remark}

%----------------------------------------------------------------------------------------------------
\subsection{The bordered triangular system}\label{bpI:p1:sec-border}

Notation: for $T:T_r\to\C^{n\times n}$ and $Y:T_r\to\C^{1\times n}$ holomorphic and $1$-periodic put
$\mathcal T:=\begin{pmatrix}\avg T&\omega\\ \avg Y&0\end{pmatrix}\in\C^{(n+1)\times(n+1)}$; with $T=T_{H,K,\omega}$,
$Y=DH(K)N$, this is $\avg{\mathsf T_H}$ (Lemma~\ref{bpI:p1:extended-basic}(a)). Let $\mathcal R$ be the solution operator of
\Kref{bp:p2:cohomological}(a): for $g$ holomorphic, periodic with $\avg g=0$, $\mathcal Rg$ is the unique zero-average
periodic solution of $\partial_\omega u=g$.

\begin{lemma}[Bordered triangular equations (HL19 Lemma 6.1)]\label{bpI:p1:border-solve}
Let $\gamma>0$, $\tau\ge0$, $\omega\in\R^n\cap\Dio_{\gamma,\tau}$, $r>0$, $\delta>0$ with $2\delta<r$, and $\Gamma_\delta=c_R\gamma^{-1}\delta^{-\varkappa}$.
Let $\eta^L,\eta^N:T_r\to\C^n$ be holomorphic, $1$-periodic, $\abs{\eta^L}\le H_L$ and $\abs{\eta^N}\le H_N$ on $T_r$,
with $\avg{\eta^N}=0$, and let $\eta^\omega\in\C$. Let $T:T_r\to\C^{n\times n}$ and $Y:T_r\to\C^{1\times n}$ be holomorphic and $1$-periodic with
$\abs T\le H_T$, $\abs Y\le H_Y$ on $T_{r-\delta}$. Assume $\det\mathcal T\ne0$ and $\abs{\mathcal T^{-1}}\le\sigma$ (row-sum
norm of $\C^{(n+1)\times(n+1)}$). Fix $\xi^L_0\in\C^n$ and define, in this order,
\begin{align*}
u^N&:=\mathcal R\eta^N\ \text{ on }T_{r-\delta},\\
\begin{pmatrix}\xi_0\\ \xi^\omega\end{pmatrix}&:=\mathcal T^{-1}\begin{pmatrix}\avg{\eta^L}+\avg{Tu^N}\\ \eta^\omega+\avg{Yu^N}\end{pmatrix}\in\C^{n}\times\C,\\
\xi^N&:=\xi_0-u^N,\qquad g^L:=T\xi^N+\omega\,\xi^\omega-\eta^L\ \text{ on }T_{r-\delta},\\
\xi^L&:=\xi^L_0+\mathcal Rg^L\ \text{ on }T_{r-2\delta}.
\end{align*}
Then:
\begin{enumerate}[label=(\alph*)]
\item (Well defined.) $u^N$, $\xi^N$, $g^L$ are holomorphic and $1$-periodic on $T_{r-\delta}$, $\avg{g^L}=0$, so $\xi^L$ is defined
  and holomorphic, $1$-periodic on $T_{r-2\delta}$; the averages $\avg{Tu^N}$, $\avg{Yu^N}$ exist.
\item (The system.) On $T_{r-2\delta}$,
  \[
  T\xi^N-\partial_\omega\xi^L+\omega\,\xi^\omega=\eta^L,\qquad-\partial_\omega\xi^N=\eta^N,\qquad\avg{Y\xi^N}=\eta^\omega,
  \]
  i.e. $\Lambda\xi-\partial_\omega\xi+\binom{\omega}{0}\xi^\omega=\eta$ with $\Lambda=\bigl(\begin{smallmatrix}0&T\\0&0\end{smallmatrix}\bigr)$,
  $\xi=(\xi^L,\xi^N)$, $\eta=(\eta^L,\eta^N)$.
\item (Estimates.) With $S_*:=\sigma\max\{H_L+H_T\Gamma_\delta H_N,\ \abs{\eta^\omega}+H_Y\Gamma_\delta H_N\}$:
  \begin{gather*}
  \abs{\xi_0},\ \abs{\xi^\omega}\le\abs{(\xi_0,\xi^\omega)}\le S_*,\qquad\sup_{T_{r-\delta}}\abs{\xi^N}\le S_*+\Gamma_\delta H_N,\\
  \sup_{T_{r-\delta}}\abs{g^L}\le H_T(S_*+\Gamma_\delta H_N)+\abs\omega S_*+H_L,\qquad
  \sup_{T_{r-2\delta}}\abs{\xi^L-\xi^L_0}\le\Gamma_\delta\sup_{T_{r-\delta}}\abs{g^L}.
  \end{gather*}
\item (Uniqueness.) If $(\tilde\xi^L,\tilde\xi^N,\tilde\xi^\omega)$ are holomorphic, periodic on $T_{r-2\delta}$ ($\tilde\xi^\omega\in\C$) and
  satisfy the three equations of (b) on $T_{r-2\delta}$, then $\tilde\xi^\omega=\xi^\omega$, $\tilde\xi^N=\xi^N$ and
  $\tilde\xi^L-\xi^L$ is constant. Thus the solution is unique up to $\avg{\xi^L}$ (the free phase direction, HL19 Rem.~3.5).
\item (Realness.) If $\omega$ is real and $\eta^L,\eta^N,T,Y$ are real on reals, $\eta^\omega\in\R$, $\mathcal T$ real (automatic), and
  $\xi^L_0\in\R^n$, then $\xi_0$, $\xi^\omega$, $\xi^N$, $\xi^L$, $g^L$ are real on reals.
\end{enumerate}
\source{HL19 Lemma~6.1, pp.~46--47, (6.7)--(6.10), with signs converted by Remark~\ref{bpI:p1:signs}; the estimates are those on
p.~46, with the additional term $\abs\omega S_*$ in $g^L$ (see Remark~\ref{bpI:p1:border-remarks}). Structure as in
\Kref{bp:p2:triangular} (HL19 Lemma~5.1, pp.~31--32).}
\sourcegap{HL19 prove Lemma~6.1 in four lines (``It is analogous to Lemma~5.1\dots''), and Lemma~5.1 is itself only referred to
(``directly adapted from Lemma~4.14 in [33]''). The proof below is supplied here in full; the new points relative to
\Kref{bp:p2:triangular} are Steps~\ref{stepI:p1:bs-bordered}, \ref{stepI:p1:bs-energy} and \ref{stepI:p1:bs-second}, and
uniqueness (d).}
\end{lemma}

\begin{proof}
\proofstep{First cohomological equation}\label{stepI:p1:bs-first}
$\eta^N$ is holomorphic, periodic on $T_r$, bounded by $H_N$, with $\avg{\eta^N}=0$. \Kref{bp:p2:cohomological}(a) gives $u^N=\mathcal R\eta^N$ holomorphic
and $1$-periodic on $T_{r-\delta}$, $\avg{u^N}=0$, $\partial_\omega u^N=\eta^N$, $\abs{u^N}\le\Gamma_\delta H_N$ on $T_{r-\delta}$, real on reals if $\eta^N$ is.

\proofstep{Averages of the products}\label{stepI:p1:bs-averages}
$Tu^N$ and $Yu^N$ are holomorphic and periodic on $T_{r-\delta}$, hence continuous on $\R^n$, so their averages exist
(\Kref{bp:p2:averages}(a)) and $\abs{\avg{Tu^N}}\le H_T\Gamma_\delta H_N$, $\abs{\avg{Yu^N}}\le H_Y\Gamma_\delta H_N$,
$\abs{\avg{\eta^L}}\le H_L$. Let $w=(\avg{\eta^L}+\avg{Tu^N},\ \eta^\omega+\avg{Yu^N})\in\C^{n+1}$; in the max norm of $\C^{n+1}$,
$\abs w\le\max\{H_L+H_T\Gamma_\delta H_N,\abs{\eta^\omega}+H_Y\Gamma_\delta H_N\}$.

\proofstep{The bordered average system}\label{stepI:p1:bs-bordered}
\emph{(new relative to \Kref{bp:p2:triangular}.)} Since $\det\mathcal T\ne0$, $(\xi_0,\xi^\omega):=\mathcal T^{-1}w$ is the unique solution of
\begin{equation}\label{bpI:p1:bordered-avg}
\avg T\xi_0+\omega\,\xi^\omega=\avg{\eta^L}+\avg{Tu^N},\qquad\avg Y\xi_0=\eta^\omega+\avg{Yu^N},
\end{equation}
and $\abs{(\xi_0,\xi^\omega)}\le\abs{\mathcal T^{-1}}\abs w\le S_*$ (\Kref{bp:p2:matrix}(b)). This is (6.10) of HL19 in our sign convention
(Remark~\ref{bpI:p1:signs}). Note that $\avg T$ alone need not be invertible: the matrix $\mathcal T$ determines the constant $\xi_0$ and the
frequency correction $\xi^\omega$ together.

\proofstep{The energy equation and $\xi^N$}\label{stepI:p1:bs-energy}
By Step~\ref{stepI:p1:bs-first}, $\xi^N=\xi_0-u^N$ is holomorphic, periodic on $T_{r-\delta}$ with $\partial_\omega\xi^N=-\eta^N$ and
$\abs{\xi^N}\le\abs{\xi_0}+\Gamma_\delta H_N\le S_*+\Gamma_\delta H_N$. By linearity of the average and the second equation of~\eqref{bpI:p1:bordered-avg},
$\avg{Y\xi^N}=\avg Y\xi_0-\avg{Yu^N}=\eta^\omega$. This is the third equation of (b) and the second one.

\proofstep{The second cohomological equation}\label{stepI:p1:bs-second}
\emph{(new relative to \Kref{bp:p2:triangular}.)} $g^L=T\xi^N+\omega\xi^\omega-\eta^L$ is holomorphic and periodic on $T_{r-\delta}$ with
$\abs{g^L}\le H_T(S_*+\Gamma_\delta H_N)+\abs\omega S_*+H_L$. Its average is, by linearity,
$\avg T\xi_0-\avg{Tu^N}+\omega\xi^\omega-\avg{\eta^L}=0$ by the first equation of~\eqref{bpI:p1:bordered-avg}. This is the point
of the border: the frequency correction $\xi^\omega$ absorbs the part of the mean of $T\xi^N-\eta^L$ that the constant $\xi_0$ cannot
adjust once it is spent on the energy equation. \Kref{bp:p2:cohomological}(a) on $T_{r-\delta}$ with loss $\delta$ (valid since $\delta<r-\delta$) gives
$\mathcal Rg^L$ holomorphic, periodic on $T_{r-2\delta}$, with $\partial_\omega\mathcal Rg^L=g^L$ and $\abs{\mathcal Rg^L}\le\Gamma_\delta\sup\abs{g^L}$. Then
$T\xi^N-\partial_\omega\xi^L+\omega\xi^\omega=T\xi^N-g^L+\omega\xi^\omega=\eta^L$. With Steps~\ref{stepI:p1:bs-first} and~\ref{stepI:p1:bs-energy} this is (a), (b), (c).

\proofstep{Uniqueness}\label{stepI:p1:bs-unique}
Let $(\tilde\xi^L,\tilde\xi^N,\tilde\xi^\omega)$ be as in (d). Then $\partial_\omega(\tilde\xi^N+u^N)=-\eta^N+\eta^N=0$ on $T_{r-2\delta}$, and
$\tilde\xi^N+u^N-\tilde c$, $\tilde c:=\avg{\tilde\xi^N+u^N}$, is a zero-average holomorphic periodic solution of $Du\,\omega=0$; by the
uniqueness in \Kref{bp:p1:cohomological} (as in Step~\ref{stepI:p1:ec-unique}) it is $0$. So $\tilde\xi^N=\tilde c-u^N$. Averaging the first equation of (b) for the tilde-solution
($\avg{\partial_\omega\tilde\xi^L}=0$, \Kref{bp:p2:averages}(b)) and using the third gives exactly~\eqref{bpI:p1:bordered-avg} for $(\tilde c,\tilde\xi^\omega)$, whence
$(\tilde c,\tilde\xi^\omega)=(\xi_0,\xi^\omega)$ by invertibility of $\mathcal T$, and $\tilde\xi^N=\xi^N$. Then $\partial_\omega(\tilde\xi^L-\xi^L)=0$, so
$\tilde\xi^L-\xi^L$ is constant, by the same uniqueness lemma applied to $\tilde\xi^L-\xi^L-\avg{\tilde\xi^L-\xi^L}$.

\proofstep{Realness}\label{stepI:p1:bs-real}
If the data are real on reals, $\avg T,\avg Y,\avg{\eta^L},\avg{Tu^N},\avg{Yu^N}$ are real (\Kref{bp:p2:averages}(a)), $\mathcal T$ is a real
matrix with real inverse (\Kref{bp:p1:real-inverse}), so $\xi_0,\xi^\omega$ are real; $\mathcal R$ preserves realness
(\Kref{bp:p2:cohomological}(a)), so $u^N,\xi^N,g^L,\xi^L$ are real on reals.
\end{proof}

\begin{remark}[Differences from HL19's printed statement]\label{bpI:p1:border-remarks}
(i) HL19's $\mathfrak R_\omega$ discards the average of its argument, so their (6.9) $\xi^L=\xi^L_0+\mathfrak R_\omega(\eta^L-T\xi^N)$ omits the
constant vector $\omega\xi^\omega$, which has no effect on non-zero Fourier modes; in the operator $\mathcal R$ of
\Kref{bp:p2:cohomological} (defined on zero-average data) the term is carried explicitly in $g^L$, which costs the extra summand
$\abs\omega S_*$ in the bound of $\sup\abs{g^L}$. HL19's bound for $\xi^L$ (p.~46, last line) is true but is not obtained from \Kref{bp:p2:cohomological}
as stated; we do not need it, because $\abs\omega\le\sigma_\omega\abs{\omega^*}$ is available in the application. (ii) HL19 have real data
and $\xi^L_0\in\R^n$; our lemma holds for complex data and records realness separately. (iii) The estimates of HL19 are in $\max$ of the two
$\C^{n+1}$ block components, as here. (iv) Uniqueness (d) is not stated in HL19.
\end{remark}

\begin{lemma}[The linearized isoenergetic equation in the frame]\label{bpI:p1:reduced}
Assume~\ref{bpI:p1:setting} with $\det G_K\ne0$ on $T_r$ and $\omega\in\R^n$ (any). Let $\xi=(\xi^L,\xi^N)$ be holomorphic on $T_{r'}$, $0<r'\le r$,
and $\xi^\omega\in\C$; put $\Delta K=P\xi$, $\Delta\omega=-\omega\xi^\omega$. Then on $T_{r'}$:
\begin{enumerate}[label=(\alph*)]
\item $QL=\binom{I_n}{-\Omega_K}$, hence $QL\omega\,\xi^\omega=\binom{\omega}{-\Omega_K\omega}\xi^\omega$.
\item (Projected identity, unconditional.) $Q\bigl(A\,\Delta K-\partial_\omega\Delta K-DK\,\Delta\omega\bigr)+QE=\Lambda\xi-\partial_\omega\xi+\binom{\omega}{0}\xi^\omega-\eta+E_{\mathrm{lin}}$, where
  \[
  E_{\mathrm{lin}}:=E_{\mathrm{red}}\xi+\Omega_0E_{\mathrm{sym}}\,\partial_\omega\xi-\binom{0}{\Omega_K\omega}\xi^\omega .
  \]
\item (Residual.) Put $\mathrm{Res}:=A\Delta K-\partial_\omega\Delta K+L\omega\xi^\omega+E$ (the residual of the linearized equation). If the bordered system
$\Lambda\xi-\partial_\omega\xi+\binom\omega0\xi^\omega=\eta$ of Lemma~\ref{bpI:p1:border-solve}(b) holds on $T_{r'}$, then $Q\,\mathrm{Res}=E_{\mathrm{lin}}$ there, and
$\abs{E_{\mathrm{lin}}}\le\abs{E_{\mathrm{red}}}\abs\xi+\abs{E_{\mathrm{sym}}}\abs{\partial_\omega\xi}+\abs{\Omega_K}\abs\omega\abs{\xi^\omega}$.
\item (Conditional equivalence.) If moreover $\det P\ne0$ on $T_{r'}$ (equivalently $\det Q\ne0$; this holds when $\abs{E_{\mathrm{sym}}}<1$, see the proof, and must be
proved in the application from its smallness conditions: it does not follow from $\det G_K\ne0$), then $\mathrm{Res}=Q^{-1}E_{\mathrm{lin}}$; and for any $\xi,\xi^\omega$,
$\mathrm{Res}=0\iff\Lambda\xi-\partial_\omega\xi+\binom\omega0\xi^\omega=\eta-E_{\mathrm{lin}}$.
\end{enumerate}
The data are $\eta^L,\eta^N$ of \Kref{bp:p2:objects}, $T$ the torsion and $Y=DH(K)N$. The defects are quadratic in the total error only on shrunken strips where the bounds of
Lemma~\ref{bpI:p1:geometry-noA5} and of Lemma~\ref{bpI:p1:border-solve}(c) hold.
\source{HL19 (6.4)--(6.5), p.~45, with $\Omega=\Omega_0$, $A=0$ (canonical Case~III); (6.34), p.~49 for $E_{\mathrm{lin}}$. Frame identities:
\Kref{bp:p2:frame-identities}.}
\end{lemma}

\begin{proof}
\proofstep{The product $QL$}\label{stepI:p1:rd-ql}
$Q=-\Omega_0P^\top\Omega_0$ and $P^\top\Omega_0L=\binom{L^\top\Omega_0L}{N^\top\Omega_0L}=\binom{\Omega_K}{I_n}$ by \Kref{bp:p2:frame-identities}(a). By
\Kref{bp:p2:matrix}(d), $\Omega_0\binom{U_1}{U_2}=\binom{-U_2}{U_1}$, so $QL=-\binom{-I_n}{\Omega_K}=\binom{I_n}{-\Omega_K}$. This is (a).

\proofstep{Expanding the linearized operator}\label{stepI:p1:rd-expand}
$\partial_\omega(P\xi)=(\partial_\omega P)\xi+P\partial_\omega\xi$, so $A\Delta K-\partial_\omega\Delta K=(AP-\partial_\omega P)\xi-P\partial_\omega\xi$, and
$-DK\,\Delta\omega=L\omega\,\xi^\omega$. By \Kref{bp:p2:frame-identities}(e), $Q(AP-\partial_\omega P)=\Lambda+E_{\mathrm{red}}$; by (f),
$QP=I_{2n}-\Omega_0E_{\mathrm{sym}}$; by (a), $QL\omega\xi^\omega=\binom\omega0\xi^\omega-\binom0{\Omega_K\omega}\xi^\omega$. Adding $QE=-\eta$:
$Q(\dots)+QE=(\Lambda+E_{\mathrm{red}})\xi-\partial_\omega\xi+\Omega_0E_{\mathrm{sym}}\partial_\omega\xi+\binom\omega0\xi^\omega-\binom0{\Omega_K\omega}\xi^\omega-\eta$, which is (b).

\proofstep{Consequence}\label{stepI:p1:rd-consequence}
If the border system holds, (b) gives $Q\,\mathrm{Res}=E_{\mathrm{lin}}$ (c). For (d): $Q=-\Omega_0P^\top\Omega_0$ is
invertible iff $P$ is, and $\det(P)^2\det\Omega_0=\det(P^\top\Omega_0P)=\det(\Omega_0+E_{\mathrm{sym}})\ne0$ when $\abs{E_{\mathrm{sym}}}<1$
(Neumann series \Kref{bp:p1:neumann} for $I-\Omega_0E_{\mathrm{sym}}$, $\abs{\Omega_0E_{\mathrm{sym}}}=\abs{E_{\mathrm{sym}}}$, \Kref{bp:p1:matrix-norms}). With $Q$ invertible,
(b) reads $Q\,\mathrm{Res}=[\Lambda\xi-\partial_\omega\xi+\binom\omega0\xi^\omega-\eta]+E_{\mathrm{lin}}$ for every $\xi,\xi^\omega$, which gives both claims of (d). Without invertibility of $Q$
multiplying by $Q$ loses information (a primary real torus with $\det G_K\ne0$ and singular $P$, $Q$ at a real angle exists). The bound on $E_{\mathrm{lin}}$ is \Kref{bp:p2:matrix}(b), (c).
\end{proof}

%----------------------------------------------------------------------------------------------------
\subsection{Geometry and frame perturbation without the torsion hypothesis}\label{bpI:p1:sec-geometry}

Track~K's Lemmas~\Kref{bp:p2:geometry} and~\Kref{bp:p2:perturbation} are stated for $(r,\beta,\varepsilon)$-admissible tori, which
include \Kref{bp:p2:A5}: invertibility of the \emph{unbordered} averaged torsion $\avg T$. In the isoenergetic setting only the
\emph{bordered} matrix $\avg{\mathsf T_H}$ is assumed invertible; $\avg T$ may be singular. The two lemmas below restate the
conclusions of K's lemmas that do not involve $\avg T$, under explicit hypotheses without any torsion condition, and prove them from
the cited K tools (the torsion hypothesis of K is used there only for its item~(f)). They are cited by Part~\ref{bpI:p2} and
\emph{replace} the K labels everywhere in the isoenergetic scheme. The regularity facts of \Kref{bp:p2:frame-basic}(a)--(d) used in the
proofs need only \Kref{bp:p2:A2}--\Kref{bp:p2:A4} (and \Kref{bp:p2:A6} for (c)); no condition on any torsion average occurs there.

\begin{lemma}[Geometry of an approximately invariant torus, without \Kref{bp:p2:A5}]\label{bpI:p1:geometry-noA5}
Let $\omega\in\Dio_{\gamma,\tau}$, $0<\delta\le1$, $2\delta<r$, $0\le\varepsilon\le1$, and let $K$ be holomorphic, primary and real on reals on $T_r$ with $B_\beta(K(\theta))\subset D$, $\abs{DK}\le\sigma_K$, $\det G_K\ne0$, $\abs B\le\sigma_B$ and $\abs{E_{H,K,\omega}}\le\varepsilon$ on $T_r$
(for some $\beta>0$; no condition on any torsion average, bordered or not; $\mathfrak m,\Theta,C_*$ are the constants of \Kref{bp:p2:constants}, depending on $(n,\tau,M,\sigma_K,\sigma_B,\sigma_T)$, with $\sigma_T$ entering only through $\mathfrak m=\max\{1,M,\sigma_K,\sigma_B,\sigma_T\}$). Then the conclusions of \Kref{bp:p2:geometry} hold:
$\abs{DE}\le n\varepsilon/\delta$, $\abs{\partial_\omega\Omega_K}\le 2n(n+1)\mathfrak m\varepsilon/\delta$, $\abs{E_{\mathrm{red}}}\le C_{\mathrm{red}}\varepsilon/\delta$ on $T_{r-\delta}$, and
$\abs{\Omega_K}\le C_\Omega\varepsilon/(\gamma\delta^{\varkappa+1})$, $\abs{E_{\mathrm{sym}}}\le C_{\mathrm{sym}}\varepsilon/(\gamma\delta^{\varkappa+1})$ on $T_{r-2\delta}$.
Also, on $T_r$: $\abs{\eta^L}\le2n\mathfrak m^2\varepsilon$, $\abs{\eta^N}\le2n\mathfrak m\varepsilon$ and $\abs L\le\mathfrak m$, $\abs B\le\mathfrak m$, $\abs N\le\mathfrak m^2$, $\abs P\le2\mathfrak m^2$; and if $\abs\omega\le\Theta$, then $\abs{X_N}\le C_X/\delta$, $\abs T\le C_T/\delta$ on $T_{r-\delta}$.
\source{HL19 Lemmas~4.3, 4.6, 4.10, pp.~22--30, as in
\Kref{bp:p2:geometry};
ours (restated without the torsion hypothesis).}
\end{lemma}
\begin{proof}
\sloppy The last sentence is \Kref{bp:p2:frame-basic}(b)--(d), whose hypotheses are \Kref{bp:p2:A2}--\Kref{bp:p2:A4} (and \Kref{bp:p2:A6} for (c)); they contain no torsion condition.
\proofstep{Derivatives of the error}\label{stepI:p1:geo-de}\sloppy
$E$ is holomorphic with $\abs E\le\varepsilon$ on $T_r$ (\Kref{bp:p2:frame-basic}(a)); \Kref{bp:p1:cauchy-strip} gives $\abs{DE}\le n\varepsilon/\delta$ on $T_{r-\delta}$.
\proofstep{$\partial_\omega\Omega_K$ and $E_{\mathrm{red}}$}\label{stepI:p1:geo-red}
By \Kref{bp:p2:frame-identities}(d), (e), \Kref{bp:p2:matrix} and \Kref{bp:p2:frame-basic}(b), as in \Kref{step:p2:geo-domegak}, \Kref{step:p2:geo-ered}:
$\abs{\partial_\omega\Omega_K}\le(2n\mathfrak m+2n^2\mathfrak m)\varepsilon/\delta$ and $\abs{E_{\mathrm{red}}}\le C_{\mathrm{red}}\varepsilon/\delta$ on $T_{r-\delta}$.
\proofstep{Lagrangian defect and symplecticity defect}\label{stepI:p1:geo-sym}
$\Omega_K$ is holomorphic, $1$-periodic on $T_{r-\delta}$ with $\avg{\Omega_K}=0$ (\Kref{bp:p2:exactness}; $K$ primary). \Kref{bp:p2:cohomological}(b) on $T_{r-\delta}$ with loss $\delta<r-\delta$ gives the bound of $\Omega_K$ on $T_{r-2\delta}$, and $\abs{E_{\mathrm{sym}}}\le\mathfrak m^2\abs{\Omega_K}$ by \Kref{bp:p2:frame-identities}(b).
\end{proof}


\begin{lemma}[Perturbation of the frame, without \Kref{bp:p2:A5}]\label{bpI:p1:perturbation-noA5}
Let $r>0$, $0<\delta\le1$, $2\delta<r$, $\abs\omega\le\Theta$, $\omega\in\R^n$. Let $K$ be holomorphic on $T_r$ with $B_\beta(K(\theta))\subset D$, $\abs{DK}\le\sigma_K$, $\det G\ne0$, $\abs B\le\sigma_B$ on $T_r$.
Let $\tilde K$ be holomorphic on $T_r$ with $\abs{\tilde K-K}\le d$ on $T_r$, where $d<\beta$ and $2\mathfrak mC_{\Delta T}\,d/\delta^2\le1$. Then (a)--(e) of \Kref{bp:p2:perturbation} hold:
\begin{enumerate}[label=(\alph*)]
\item $B_{\beta-d}(\tilde K(\theta))\subset D$ on $T_r$ and $\abs{\tilde A-A}\le\mathfrak md$;
\item $\abs{D\tilde K-DK}\le nd/\delta$ on $T_{r-\delta}$;
\item $\det\tilde G\ne0$, $\abs{\tilde B}\le2\mathfrak m$, $\abs{\tilde B-B}\le C_Bd/\delta$ on $T_{r-\delta}$;
\item $\abs{\tilde N-N}\le C_{\Delta N}d/\delta$, $\abs{\tilde N}\le2\mathfrak m^2$, $\abs{\tilde N^\top}\le4n\mathfrak m^2$ on $T_{r-\delta}$;
\item $\abs{\tilde T-T}\le C_{\Delta T}d/\delta^2$ on $T_{r-2\delta}$ (here $\tilde T$ is the torsion of $\tilde K$ for the same $\omega$).
\end{enumerate}
\source{HL19 Lemma~5.2, Steps~3--7, pp.~36--42, as in \Kref{bp:p2:perturbation}; the torsion hypothesis there is used only for its item (f).}
\end{lemma}
\begin{proof}
\proofstep{Smallness}\label{stepI:p1:pt-small}
From the definitions of \Kref{bp:p2:constants}, $C_{\Delta T}\ge\max\{n,C_G,C_{\Delta N}\}$ (\Kref{step:p2:pt-small}, pure arithmetic on constants). With $d/\delta\le d/\delta^2\le1/(2\mathfrak mC_{\Delta T})$:
$nd/\delta\le\mathfrak m$, $\mathfrak mC_Gd/\delta\le\frac12$, $C_{\Delta N}d/\delta\le\mathfrak m^2$, $\mathfrak mC_{\Delta T}d/\delta^2\le\frac12$.
\proofstep{Domain}\label{stepI:p1:pt-domain}
$\abs{z-\tilde K(\theta)}<\beta-d$ implies $\abs{z-K(\theta)}<\beta$, so $z\in D$; the segment $[K,\tilde K]$ lies in $B_\beta(K(\theta))\subset D$, and \Kref{bp:p2:hamiltonian}(e) gives $\abs{\tilde A-A}\le Md\le\mathfrak md$.
\proofstep{Frame}\label{stepI:p1:pt-frame}
\Kref{bp:p1:cauchy-strip}, \Kref{bp:p1:cauchy-corollaries} ($p=2n$) on $\tilde K-K$: $\abs{\tilde L-L}\le nd/\delta$, $\abs{(\tilde L-L)^\top}\le2nd/\delta$ on $T_{r-\delta}$; so $\abs{\tilde L}\le2\mathfrak m$, $\abs{\tilde L^\top}\le4n\mathfrak m$. Then $\tilde G-G=\tilde L^\top(\tilde L-L)+(\tilde L-L)^\top L$ has norm $\le C_Gd/\delta$, $\abs B\abs{\tilde G-G}\le\frac12$, and \Kref{bp:p1:neumann} gives $\tilde G$ invertible, $\abs{\tilde B}\le2\mathfrak m$, $\abs{\tilde B-B}\le C_Bd/\delta$ (Gram invertibility before any use of $\tilde N$, $\tilde T$: Remark~\ref{bpI:p0:totality}(T2)).
\proofstep{Normal frame and torsion}\label{stepI:p1:pt-torsion}
$\tilde N-N=\Omega_0\bigl((\tilde L-L)\tilde B+L(\tilde B-B)\bigr)$ gives (d). On $T_{r-2\delta}$,
$\tilde X_{\tilde N}-X_N=(\tilde A-A)\tilde N+A(\tilde N-N)-\partial_\omega(\tilde N-N)$, with the first two terms $\le2\mathfrak m^3d+\mathfrak mC_{\Delta N}d/\delta$ and the third $\le n\Theta C_{\Delta N}d/\delta^2$ (\Kref{bp:p1:cauchy-corollaries} with $\abs\omega_1\le n\Theta$); so $\abs{\tilde X_{\tilde N}-X_N}\le C_{\Delta X}d/\delta^2$. Finally
$\tilde T-T=\tilde N^\top\Omega_0(\tilde X_{\tilde N}-X_N)+(\tilde N-N)^\top\Omega_0X_N$ and $\abs{X_N}\le C_X/\delta$ on $T_{r-\delta}$ (\Kref{bp:p2:frame-basic}(d), whose hypotheses \Kref{bp:p2:A2}--\Kref{bp:p2:A4} and $\abs\omega\le\Theta$ are exactly those of this lemma; no error or torsion condition) give $C_{\Delta T}d/\delta^2$.
\end{proof}


%----------------------------------------------------------------------------------------------------
\subsection{Suspected errata and sign checks}\label{bpI:p1:sec-errata}

\begin{remark}[Candidate errata of HL19, checked against the pinned PDF {\upshape(\texttt{arxiv\_1711.05515.pdf}, v1)}]\label{bpI:p1:errata}
All three were checked on the rendered pages 18, 19, 46 (not only on the text layer).
\begin{description}
\item[E1 (strict inequalities).] Thm~3.6, (3.16), p.~19: ``$\|K_\infty-K\|_{\rho_\infty}<\dfrac{\mathfrak C_2\|E_c\|_\rho}{\gamma^2\rho^{2\tau}},\quad
  |\omega_\infty-\omega|<\dfrac{\mathfrak C_3\|E_c\|_\rho}{\gamma^2\rho^{2\tau}}$''; likewise (3.11), p.~18 for Thm~3.1. With $E_c\equiv0$ the right side is $0$,
  and (3.15) holds, but no $K_\infty$ satisfies $\|K_\infty-K\|<0$. The statements must read ``$\le$'' (the frozen I-G1 does).
\item[E2 (sign in (3.12)).] p.~18: ``(3.12) $X_h(K(\theta))=\mathfrak L_\omega K(\theta)$, $\langle c\circ K\rangle=c_0$''. With
  $\mathfrak L_\omega=-\partial_\omega$ (4.1), p.~20, this reads $X_h(K)=-DK\omega$, contradicting the invariance equation
  (2.1), p.~6, $X_h(K)=DK\,\omega$, and the error (2.2), p.~7, $E=X_h(K)-DK\omega$. The correct form is $X_h(K)+\mathfrak L_\omega K=0$, which is how
  p.~45 (the total error before (6.1)) and p.~48 (Lemma~6.2, $\bar E_c$) write it. Our Remark~\ref{bpI:p1:signs} uses the consistent convention.
\item[E3 (Rem.~3.7, first equality).] p.~19: ``$Dh(K(\theta))N(\theta)=X_h(K(\theta))^\top\Omega(K(\theta))N(\theta)=-\omega^\top DK(\theta)^\top\Omega(K(\theta))N(\theta)=(\omega^\top\ 0^\top_{n-d})=\hat\omega^\top$''.
  By Rem.~2.1, p.~6, $\Omega X_h=(Dh)^\top$, so $Dh=(\Omega X_h)^\top=-X_h^\top\Omega$: the first equality misses a minus sign. The second equality is
  likewise wrong for $X_h=+DK\omega$ (it would need $X_h=-DK\omega$, the E2 convention). The two slips cancel and the final value
  $Dh\,N=\hat\omega^\top$ is correct, as Lemma~\ref{bpI:p1:last-row}(a) shows (using $L^\top\Omega_0N=-I_n$, \Kref{bp:p2:frame-identities}(a)); an independent check at the trivial torus is
  Lemma~\ref{bpI:p1:trivial}, Step~\ref{stepI:p1:tr-border}. The correct chain is
  $Dh\,N=-X_h^\top\Omega N=-\omega^\top DK^\top\Omega N=\omega^\top$.
\end{description}
Other points checked and found consistent (no erratum): (6.7)--(6.10), p.~46, with $\mathfrak L_\omega=-\partial_\omega$, $\mathfrak R_\omega=-\mathcal R$
(Remark~\ref{bpI:p1:signs}); (6.3) and $\bar\omega=(1-\xi^\omega)\omega$ (p.~50, (6.32)); the equation $QL\omega=\binom\omega0-\binom0{\Omega_K\omega}$ against the display on p.~46
(Lemma~\ref{bpI:p1:reduced}(a)); the sign $\Delta\omega=-\omega\xi^\omega$. HL19's (6.9) omits the constant $\omega\xi^\omega$ (Remark~\ref{bpI:p1:border-remarks}(i)), which is harmless.
\end{remark}

\begin{remark}[Doubts about the frozen statements]\label{bpI:p1:doubts}
None of the frozen statements used in this part looks false. Two reading points. (a) The frozen $\mathsf B_h$ uses the real derivative of $\operatorname{Re}h$ via
\Lean{frequencyMap}; Lemma~\ref{bpI:p1:trivial} identifies it with the complex Hessian of $h$ only under ``$h$ real on the real points of a neighbourhood of $I$'',
which I-HEAD assumes on all of $B_{r_0}$; without that hypothesis $\mathsf T_h$ and $\mathsf B_h$ can differ by the imaginary part. (b) The frozen last row of $\mathsf T_H$ is
$DH(K)N$, not $\omega^\top$ (Definition~\ref{bpI:p0:extended}); Lemma~\ref{bpI:p1:last-row} shows the difference is $(\eta^L)^\top=O(e)$, so the bordered matrices with the frozen row and with the invariant row are perturbatively comparable once a Neumann bound
(\Kref{bp:p1:neumann}) and the associated smallness and margin conditions are proved; this is not automatic, because $\sigma_T$ is fixed before the error in I-G1. The frozen row $DH(K)N$ is retained throughout.
\end{remark}
